% =====================================================================
\section{The Quadcopter Model: Blocks 7--10}
\label{sec:x3}
% =====================================================================

The paper controls a body that can produce torque about every axis directly and does not need to hold itself up. A
quadcopter is different in three ways: (i) its four rotors produce one thrust along the body $z$-axis plus three torques,
so attitude and position are coupled (it moves sideways by tilting); (ii) Gazebo publishes only the \emph{pose}, not the
angular velocity; (iii) the controller runs as a sampled loop at the pose rate ($\approx$60\,Hz) with a transport delay.
Blocks 7--10 are what we added so that blocks 3--6 can be used unchanged on the X3. First, the physics that Gazebo
integrates (and that our software-in-the-loop surrogate \code{x3\_sim/x3\_model.py} reproduces).

\begin{eqbox}{cPlant}{The X3 quadcopter as Gazebo simulates it}
\begin{align}
m\ddot{\vp}&=R\begin{bmatrix}0\\0\\\sum_iF_i\end{bmatrix}-mg\,\bm e_3-\Big(\sum_i|\Omega_i|\Big)c_D\,\vv_\perp,\qquad F_i=k_T\Omega_i^2 \tag{X1}\label{eq:x3trans}\\
\dot\vq&=\tfrac12\,\vq\otimes[\vw;0],\qquad J\dot\vw=-\vw^\times J\vw+\vtau,\qquad
\vtau=\sum_i\Big(\vr_i\times F_i\bm e_3-\mathrm{dir}_i\,c_QF_i\,\bm e_3\Big) \tag{X2}\label{eq:x3rot}\\
\dot\Omega_i&=\frac{\Omega_{i,cmd}-\Omega_i}{\tau_{\uparrow/\downarrow}},\qquad \tau_\uparrow=12.5\ \text{ms},\ \tau_\downarrow=25\ \text{ms},\qquad 0\le\Omega_{i,cmd}\le800\ \text{rad/s} \tag{X3}\label{eq:x3rotor}
\end{align}
\end{eqbox}

\begin{dimtable}
\vp,\ \vv & position and velocity of the CoM (world) & \R^3 & m, m/s\\
R & rotation matrix of $\vq$ (body $\to$ world) & \R^{3\times3} & --\\
F_i,\ \Omega_i & thrust and speed of rotor $i$ & \R & N, rad/s\\
\vr_i & position of rotor $i$ in the body frame & \R^3 & m\\
\mathrm{dir}_i & $+1$ counter-clockwise, $-1$ clockwise rotor & \R & --\\
c_D & rotor drag coefficient ($8.06\times10^{-5}$) & \R & N\,s$^2$/(m\,rad)\\
\vv_\perp & velocity component perpendicular to the rotor axis & \R^3 & m/s\\
\end{dimtable}

(X2) is exactly the paper's Eq.~(1) with $\vu+\vd$ replaced by the torque of the rotors: the attitude part of the X3
\emph{is} the paper's plant, with a 250 times smaller inertia. (X1) is what the paper does not have.

% ---------------------------------------------------------------------
\subsection{Block 7 --- Body rate from two poses}
\label{sec:b7}
% ---------------------------------------------------------------------

\begin{idea}
If we know the orientation now and one sample ago, the small rotation between them, divided by the time between
them, is the body rate. The small rotation must be computed \emph{in the body frame} ($\vq_{k-1}^*\otimes\vq_k$), not by
subtracting quaternion components.
\end{idea}

\begin{eqbox}{cX3}{Block 7: exact body rate from consecutive quaternions}
\begin{equation}
\delta\vq_k=\vq_{k-1}^{*}\otimes\vq_k,\qquad
\vw_{raw}=\frac{2\,\mathrm{atan2}\big(\|\delta\vq_v\|,\,\delta q_4\big)}{\Delta t_k}\;\frac{\delta\vq_v}{\|\delta\vq_v\|},\qquad
\vw_k=0.4\,\vw_{k-1}+0.6\,\vw_{raw} \tag{B7}\label{eq:b7}
\end{equation}
$\Delta t_k$ is the difference of the \emph{message time stamps} (simulation time), never of the wall clock.
\end{eqbox}

\subsubsection*{Derivation}
From (1a) over one sample with constant $\vw$, $\vq_k=\vq_{k-1}\otimes\delta\vq$ with
$\delta\vq=[\bm n\sin\tfrac{\phi}{2};\cos\tfrac{\phi}{2}]$ and $\phi=\|\vw\|\Delta t$, $\bm n=\vw/\|\vw\|$. Multiplying on the left by
$\vq_{k-1}^*$ isolates $\delta\vq$; its angle is $\phi=2\,\mathrm{atan2}(\|\delta\vq_v\|,\delta q_4)$ and its axis is
$\delta\vq_v/\|\delta\vq_v\|$, so $\vw=\phi\,\bm n/\Delta t$. The first-order filter with weight 0.6 removes sampling noise
with a lag of about $\tfrac{0.4}{0.6}\Delta t\approx11$\,ms at 60\,Hz.

\begin{toybox}[why the naive formula fails]
The drone is yawed by $90^\circ$, $\vq_{k-1}=[0,0,0.7071,0.7071]$, and rolls about its own $x$-axis at 0.6\,rad/s. After
$\Delta t=1/60$\,s it has turned by $0.01$\,rad: $\vq_k=[0.00354,\ 0.00354,\ 0.70710,\ 0.70710]$.
\begin{itemize}[leftmargin=*]
  \item Exact (B7): $\delta\vq=[0.0050,0,0,1.0000]$, angle $0.0100$\,rad about $x$, so $\vw=[0.600,\ 0,\ 0]$\,rad/s.
  \item Naive $2(\vq_{k,v}-\vq_{k-1,v})/\Delta t$ (our first Gazebo node): $[0.424,\ 0.424,\ -0.001]$\,rad/s: the roll is
        reported as half roll, half pitch. The naive formula equals the body rate only when the vehicle has no yaw.
\end{itemize}
\end{toybox}

\begin{projbox}
\code{FixedController.update()} in \code{x3\_sim/controllers.py}, section ``F11''. The first node, \code{legacy/gazebo\_v1.py},
used the naive difference with a filter weight of 0.10, i.e.\ a lag of about $9\Delta t\approx150$\,ms.
\end{projbox}

% ---------------------------------------------------------------------
\subsection{Block 8 --- Outer loop: thrust and the attitude to aim for}
\label{sec:b8}
% ---------------------------------------------------------------------

\begin{idea}
To hover at a point the drone needs (a) a total thrust that holds it up and (b) a tilt that points part of the
thrust towards the point. A PD law on position gives the acceleration we want; that acceleration plus gravity gives the
direction the thrust must point in, which fixes the desired attitude. The paper's law then only has to turn the body
to this attitude: its ``error quaternion'' is the rotation from the desired to the actual attitude.
\end{idea}

\begin{eqbox}{cX3}{Block 8: position loop, thrust and desired attitude}
\begin{align}
\bm a_{cmd}&=\ddot\vp_d+K_p(\vp_d-\vp)+K_d(\vv_d-\vv)\quad(\text{horizontal part limited to a tilt of }0.35\ \text{rad}) \tag{B8a}\\
\vF_d&=m(\bm a_{cmd}+g\bm e_3),\qquad T=\vF_d^\top R\,\bm e_3 \tag{B8b}\\
\bm z_{b,d}&=\frac{\vF_d}{\|\vF_d\|},\quad \bm y_{b,d}=\frac{\bm z_{b,d}\times\bm x_c}{\|\bm z_{b,d}\times\bm x_c\|},\quad
\bm x_{b,d}=\bm y_{b,d}\times\bm z_{b,d},\quad \bm x_c=[\cos\psi_d,\sin\psi_d,0]^\top \tag{B8c}\\
\vq_e&=\vq_d^*\otimes\vq,\qquad \vw_e=\vw-R_e^\top\vw_d \tag{B8d}
\end{align}
Gains: $K_{p,z}=5$, $K_{d,z}=3.5$ (vertical, acceleration clipped to $\pm3$\,m/s$^2$), $K_{p,xy}=1.2$, $K_{d,xy}=1.6$.
\end{eqbox}

\subsubsection*{Derivation}
Projecting (X1) on the world axes and neglecting the small rotor drag, $m\ddot\vp=T\,R\bm e_3-mg\bm e_3$. To get the
acceleration $\bm a_{cmd}$ we need $T\,R\bm e_3=m(\bm a_{cmd}+g\bm e_3)=\vF_d$: the body $z$-axis must point along $\vF_d$
(B8c), and the thrust is the part of $\vF_d$ along the \emph{current} body axis (B8b), so a tilted drone does not
over-thrust. The desired yaw $\psi_d$ fixes the remaining rotation about $\bm z_{b,d}$. With $\vq_e$ as the attitude error,
$\qv\to\vq_{e,v}$ in every equation of blocks 3--6, and $\vw_e$ adds the rate of the reference as a feed-forward
(it is zero in hover).

\begin{toybox}[a 10\,cm position error in hover]
$\vp_d-\vp=[0.1,0,0]$\,m, $\vv=0$: $\bm a_{cmd}=[0.12,\ 0,\ 0]$\,m/s$^2$, so
$\vF_d=1.52\,[0.12,\ 0,\ 9.81]=[0.182,\ 0,\ 14.911]$\,N. The thrust axis must lean by
$\arctan(0.182/14.911)=0.70^\circ$ towards $+x$, i.e.\ $\vq_d=[0,\ 0.00612,\ 0,\ 0.99998]$ (a pitch of $0.70^\circ$), and
$T=14.91$\,N, which is the hover thrust $mg=1.52\times9.81=14.911$\,N.
\end{toybox}

% ---------------------------------------------------------------------
\subsection{Block 9 --- The paper's law on a 60\,Hz loop}
\label{sec:b9}
% ---------------------------------------------------------------------

\begin{idea}
The paper's law is written in continuous time: its barrier gain may grow without bound as $|\varepsilon|\to1$. A digital
loop that only sees the drone every $1/60$\,s and acts one sample late cannot use an unbounded derivative gain: past a
certain value it overshoots more every sample and oscillates. So the sum of all gains acting on $\vs$ must stay below the
largest value the sampled loop tolerates. We also re-size the envelope and the gains for a body 250 times lighter.
\end{idea}

\begin{eqbox}{cCtrl}{Block 9: blocks 3--6 on the X3}
\begin{align}
g_i&=\min\Big(\frac{\alpha}{(1-\varepsilon_{1i})^2}+\frac{\beta}{(1-\varepsilon_{2i})^2},\ \ g_{\max}\frac{J_i}{J_z}\Big),\qquad
\vu=-k_p\vq_{e,v}-k_d\vw_e-\big[g_is_i+\that\,\sat(s_i/\varsigma)\big]_i \tag{B9a}\\
\dot\that&=\frac{1}{\eta_1}\big(\|\vs\|-\max(2e^{-0.5t},0.2)\,\that\big),\qquad 0\le\that\le0.1,\qquad \vtau=\sat_{0.35}(\vu) \tag{B9b}
\end{align}
X3 values: $k_p=4$, $k_d=0.4$, $\alpha=0.15$, $\beta=0.12$, $c=1$, $\varsigma=0.5$, $\eta_1=5$, $g_{\max}=1$;
envelopes $\rho_q:0.30\to0.02$, $\rho_\omega:2.5\to0.05$\,rad/s, $\kappa=2$, $T_s=10$\,s.
\end{eqbox}

\subsubsection*{Sizing the gains}
For small angles $\vq_{e,v}\approx\tfrac12\bm\phi$, so one axis of the PD part is $J\ddot\phi=-\tfrac{k_p}{2}\phi-k_d\dot\phi$, a
mass--spring--damper with $\omega_n=\sqrt{k_p/(2J)}$ and $\zeta=k_d/(2J\omega_n)$. For the paper ($J\approx17$, $k_p=40$,
$k_d=20$): $\omega_n=1.08$\,rad/s, $\zeta=0.54$. For the X3 roll axis ($J=0.0357$, $k_p=4$, $k_d=0.4$): $\omega_n=7.5$\,rad/s,
$\zeta=0.75$ --- about seven times faster, as a small drone must be, and better damped.

\subsubsection*{The sampled-data limit}
One axis, zero-order hold $T=1/60$\,s, one sample of delay: with $\bm x_k=[\phi_k;\dot\phi_k]$,
\[
\bm x_{k+1}=\underbrace{\begin{bmatrix}1&T\\0&1\end{bmatrix}}_{A_d}\bm x_k+\underbrace{\begin{bmatrix}T^2/(2J)\\T/J\end{bmatrix}}_{B_d}u_{k-1},\qquad
u_k=-\tfrac{k_p}{2}\phi_k-k_{d,tot}\dot\phi_k .
\]
Stacking $[\bm x_k;u_{k-1}]$ gives a $3\times3$ closed-loop matrix
$\begin{bmatrix}A_d&B_d\\-\tfrac{k_p}{2}\ \ -k_{d,tot}&0\end{bmatrix}$; the loop is stable when all its eigenvalues satisfy $|\lambda|<1$.
$k_{d,tot}=k_d+g_i$ is the total gain on the rate, because the barrier multiplies $\vs=\vw+c\qv$.

\begin{toybox}[how large may the barrier gain become?]
Roll axis, $J=0.03565$: $B_d=[T^2/(2J);\ T/J]=[0.003896;\ 0.4675]$. Scanning $k_{d,tot}$, the spectral radius stays below 1
for $0.06\le k_{d,tot}\le\mathbf{2.12}$\,N\,m\,s; for pitch the limit is $4.20$ and for yaw $5.91$. With the cap, the largest
total gains are $0.4+1.0\cdot0.360=0.76$ (roll), $0.4+0.711=1.11$ (pitch) and $1.40$ (yaw): all well inside. Without
the cap, $\alpha/(1-\varepsilon)^2$ alone reaches $0.15/(1-0.99)^2=1500$ at $\varepsilon=0.99$, far beyond any limit.
\end{toybox}

\begin{codebox}
\textbf{MATLAB:} \code{toy\_model/step8\_sampled\_data\_limit.m} builds the matrix above and plots the spectral radius against
$k_{d,tot}$ for the three axes.
\end{codebox}

\begin{warn}[What the cap costs]
With the cap the barrier can no longer become infinite, so the controller on the X3 can no longer \emph{force} the state
to stay inside the tube no matter what; it can only push harder near the wall, up to the cap. Any guarantee of the
paper is a continuous-time guarantee, and this is the price of running it on a real digital loop. The complete
documentation measures how often the envelope is left in each experiment.
\end{warn}

% ---------------------------------------------------------------------
\subsection{Block 10 --- From thrust and torque to four rotor speeds}
\label{sec:b10}
% ---------------------------------------------------------------------

\begin{idea}
Each rotor pushes up with $F_i$. Pushing up at a point off the centre makes a torque: a rotor at lateral position $y_i$
gives a torque $y_iF_i$ about $x$ (roll); a rotor in front ($x_i>0$) pitches it with $-x_iF_i$; and each rotor's drag twists the body
about $z$ with $-\mathrm{dir}_ic_QF_i$. These four sums form a $4\times4$ matrix. Inverting it gives the four thrusts that
produce exactly the requested thrust and torques.
\end{idea}

\begin{eqbox}{cX3}{Block 10: allocation (mixer) from the X3 geometry}
\begin{equation}
\begin{bmatrix}T\\\tau_x\\\tau_y\\\tau_z\end{bmatrix}=A\vF,\qquad
A=\begin{bmatrix}1&1&1&1\\-0.22&0.20&0.22&-0.20\\-0.13&0.13&-0.13&0.13\\-0.016&-0.016&0.016&0.016\end{bmatrix},\qquad
\vF=A^{-1}\begin{bmatrix}T\\\vtau\end{bmatrix},\quad \Omega_i=\sqrt{F_i/k_T} \tag{B10}\label{eq:mixer}
\end{equation}
\end{eqbox}

\subsubsection*{Derivation}
$\vr_i\times F_i\bm e_3=[y_iF_i,\ -x_iF_i,\ 0]^\top$, and the yaw reaction of rotor $i$ is $-\mathrm{dir}_ic_QF_i$ (the
plugin's sign convention). Summing over the rotors gives the rows of $A$ with the rotor positions of the SDF:
rotor 0 $(0.13,-0.22)$ ccw, rotor 1 $(-0.13,0.20)$ ccw, rotor 2 $(0.13,0.22)$ cw, rotor 3 $(-0.13,-0.20)$ cw. $A$ is
invertible, so any $(T,\vtau)$ maps to unique thrusts; physically they must also satisfy $0\le F_i\le k_T\cdot800^2=5.471$\,N.
If not, our mixer first drops yaw, then scales roll and pitch down, so that thrust (staying in the air) has priority.

\begin{toybox}[hover and the torque the X3 can produce]
Hover: $T=mg=14.911$\,N, $\vtau=0\Rightarrow F_i=3.728$\,N each and $\Omega_i=\sqrt{3.728/8.54858\times10^{-6}}=660.4$\,rad/s.

Largest torque about one axis while hovering (each $F_i$ must stay in $[0,\,5.471]$): roll $1.464$, pitch $0.907$, yaw
$\mathbf{0.106}$\,N\,m. The paper's disturbance has a yaw amplitude of $0.15$\,N\,m, more than the X3 can ever cancel: at
full paper amplitude the yaw constraint is infeasible for \emph{any} controller.
\end{toybox}

\begin{toybox}[what our first mixer did]
The first node used a textbook ``+/$-$'' mixer, $m_i=\tfrac{T}{4k_T}\pm\tfrac{\tau_y}{4k_Tl}\mp\tfrac{\tau_x}{4k_Tl}\mp\tfrac{\tau_z}{4c_Q}$,
whose signs do not match the X3's rotor numbering. Pushing it through $A$ gives the torque produced per unit commanded:
\[
M_{v1}=\begin{bmatrix}0.044&-0.933&0\\0&0&6.9\times10^{-5}\\0.071&0&0\end{bmatrix}\qquad(\text{ideal: }I_3).
\]
A pitch command produced a \emph{roll} torque, nothing produced a pitch torque, and the yaw term (divided by $c_Q$ instead
of $k_Tc_Q$) was five orders of magnitude too small. With (B10), $M=A_{2:4}A^{-1}_{:,2:4}=I_3$ exactly.
\end{toybox}

\begin{codebox}
\textbf{MATLAB:} \code{toy\_model/step7\_x3\_mixer.m} builds $A$ from the SDF, prints the hover speeds, the authority and
$M_{v1}$, and plots $M_{v1}$ next to the corrected mixer.
\end{codebox}

\begin{projbox}
\code{allocation\_matrix()} in \code{x3\_sim/x3\_model.py} builds $A$; the controller stores $A^{-1}$ and the de-saturation
loop is at the end of \code{FixedController.update()}. The rotor speeds are published on
\code{/X3/gazebo/command/motor\_speed} by \code{gazebo.py}.
\end{projbox}
